% unidades/18-comparaciones.tex
\chapter{⚖️ 比較 — Comparaciones}
\unidadheader{18}{比較}{Comparaciones}{Comparar personas, lugares y cosas}

\begin{tcolorbox}[vocabbox, title={📌 Vocabulario Nuevo}]
\begin{tabularx}{\linewidth}{llX}
\toprule
\textbf{Japonés} & \textbf{Español} & \textbf{Tipo} \\
\midrule
\vocabitem{一番}{いちばん}{el/la más ~ (superlativo)}{adv.}
\vocabitem{〜より}{—}{más que ~ (comparativo)}{part.}
\vocabitem{〜の方が}{ほうが}{~ es mejor / más ~ que}{expr.}
\vocabitem{同じ}{おなじ}{igual / mismo}{adj-な}
\vocabitem{違う}{ちがう}{diferente / equivocado}{v. gr.1}
\vocabitem{もっと}{—}{más (cantidad/grado)}{adv.}
\vocabitem{少し}{すこし}{un poco}{adv.}
\vocabitem{ずっと}{—}{mucho más / todo el tiempo}{adv.}
\vocabitem{だいぶ}{—}{bastante / considerablemente}{adv.}
\vocabitem{比べる}{くらべる}{comparar}{v. gr.2}
\bottomrule
\end{tabularx}
\end{tcolorbox}

\subsection*{🗣️ Frases Clave}

\frase{AとBとどちらが\ruby{好}{す}きですか？}{¿Cuál prefieres, A o B?}
\frase{Aの\ruby{方}{ほう}が\ruby{好}{す}きです.}{Prefiero A.}
\frase{\ruby{日本}{にほん}で\ruby{一番}{いちばん}\ruby{高}{たか}い\ruby{山}{やま}は？}
  {¿Cuál es la montaña más alta de Japón?}
\frase{\ruby{富士山}{ふじさん}が\ruby{一番}{いちばん}\ruby{高}{たか}いです.}{El Monte Fuji es el más alto.}

\subsection*{⚙️ Gramática}

\patron{Comparativo: más que ~}
  {A は B より 〜です}
  {\ej{\ruby{東京}{とうきょう}は\ruby{大阪}{おおさか}より\ruby{大}{おお}きいです。}
    {Tokio es más grande que Osaka.}
  \ej{\ruby{電車}{でんしゃ}の\ruby{方}{ほう}がバスより\ruby{速}{はや}いです。}
    {El tren es más rápido que el autobús.}}

\patron{Superlativo: el/la más ~}
  {〜の中で 一番 〜です}
  {\ej{このクラスの中で\ruby{誰}{だれ}が\ruby{一番}{いちばん}\ruby{背}{せ}が\ruby{高}{たか}いですか？}
    {¿Quién es el más alto de esta clase?}
  \ej{\ruby{私}{わたし}が\ruby{一番}{いちばん}\ruby{背}{せ}が\ruby{高}{たか}いです。}
    {Yo soy el más alto.}}

\begin{tcolorbox}[ejerciciobox, title={📝 Ejercicios Rápidos}]
\begin{enumerate}
  \item Compara dos ciudades que conoces usando より.
  \item ¿Cuál es tu materia favorita? Usa 一番 para expresarlo.
  \item Traduce: コーヒーとお茶とどちらの方が好きですか？
\end{enumerate}
\vspace{0.4em}
\textbf{Respuestas:}
1) (personal)\quad
2) (personal)\quad
3) ¿Cuál prefieres, el café o el té?
\end{tcolorbox}

\begin{tcolorbox}[kanjibox, title={🀄 Kanjis de la Unidad}]
\begin{tabularx}{\linewidth}{cllXX}
\toprule
\textbf{Kanji} & \textbf{On} & \textbf{Kun} & \textbf{Significado} & \textbf{Vocab} \\
\midrule
\kanjirow{一}{いち}{ひと}{uno}{\ruby{一番}{いちばん} / \ruby{一人}{ひとり}}
\kanjirow{番}{ばん}{—}{número / turno}{\ruby{一番}{いちばん} / \ruby{交番}{こうばん}}
\kanjirow{比}{ひ}{くら}{comparar}{\ruby{比較}{ひかく} / \ruby{比}{くら}べる}
\kanjirow{較}{かく}{—}{comparar}{\ruby{比較}{ひかく}}
\kanjirow{同}{どう}{おな}{igual}{\ruby{同じ}{おなじ} / \ruby{同意}{どうい}}
\bottomrule
\end{tabularx}
\end{tcolorbox}
